\documentclass[14pt,a4paper]{extarticle}
\usepackage[T2A]{fontenc}
\usepackage[utf8]{inputenc}
\usepackage[english,russian]{babel}
\usepackage{lmodern}
\usepackage{geometry}
\usepackage{array}
\usepackage{tabularx}
\usepackage{longtable}
\usepackage{ltablex}
\usepackage{booktabs}
\usepackage{hyperref}
\usepackage{indentfirst}
\usepackage{microtype}
\usepackage{graphicx}
\usepackage{setspace}
\usepackage{listings}
\usepackage{graphicx}
\usepackage{caption}
\usepackage{float}
\usepackage{tocloft}
\usepackage{titlesec}
\usepackage{newtxtext,newtxmath}

\titleformat{\section}        % форматируем \section
  {\normalfont\bfseries}      % стиль: обычный + жирный
  {\thesection.}              % номер секции с точкой
  {2mm}                       % расстояние между номером и текстом
  {}                           % перед заголовком можно добавить что-то
\titlespacing*{\section}{1.5cm}{*3}{*1}  % {отступ слева}{отступ сверху}{отступ снизу}

\titleformat{\subsection}        % форматируем \section
  {\normalfont\bfseries}      % стиль: обычный + жирный
  {\thesubsection.}              % номер секции с точкой
  {2mm}                       % расстояние между номером и текстом
  {}                           % перед заголовком можно добавить что-то
\titlespacing*{\subsection}{1.5cm}{*3}{*1}  % {отступ слева}{отступ сверху}{отступ снизу}

\captionsetup{justification=centering}
\renewcommand{\ttdefault}{cmtt}
\keepXColumns
\setlength{\parindent}{3em}
\newgeometry{left=2.5cm,right=2.5cm,top=3cm,bottom=2.5cm}

\renewcommand{\contentsname}{\hspace*{4em}{\smallСодержание}}
\setlength{\cftbeforesecskip}{0.5em}
\renewcommand{\cftsecleader}{\cftdotfill{\cftdotsep}}
\renewcommand{\cftsecaftersnum}{.\hspace{2mm}}

\title{Отчёт по "Управление ИТ-проектами"}
\author{Выполнил студент группы М-РИТ-25-1, Станиславчук Сергей}

\begin{document}

\setstretch{1.5}
\noindent
\begin{minipage}[t]{0.01\textwidth}
\includegraphics{/home/stanik/programmer/++Programmer/5_2/LSTU.jpg}
\end{minipage}%
\begin{minipage}[c]{1.1\linewidth}
    \centering
    {\small
    МИНИСТЕРСТВО НАУКИ И ВЫСШЕГО ОБРАЗОВАНИЯ РФ \\[1em]
    ФЕДЕРАЛЬНОЕ ГОСУДАРСТВЕННОЕ БЮДЖЕТНОЕ \\
    ОБРАЗОВАТЕЛЬНОЕ УЧРЕЖДЕНИЕ ВЫСШЕГО ОБРАЗОВАНИЯ \\
    «ЛИПЕЦКИЙ ГОСУДАРСТВЕННЫЙ ТЕХНИЧЕСКИЙ УНИВЕРСИТЕТ»
    }
\end{minipage}%

\vspace{3em}

\noindent
\hspace{1em}%
\begin{minipage}[t]{0.8\textwidth}
\textbf{Институт} \hspace{2em} компьютерных наук \\
\textbf{Кафедра} \hspace{2em} автоматизированных систем управления
\end{minipage}


\vspace{6em}

\begin{center}
\textbf{Практическая работа №15} \\
по дисциплине \\ <<Управление ИТ-проектами>>
\end{center}

\vspace{6em}

\noindent
\begin{tabular*}{\textwidth}{@{\extracolsep{\fill}} l c r }
    Студент М-РИТ-25-1 & \makebox[4cm]{\hrulefill} & Станиславчук С.М. \\[2em]
Руководитель & \\
К.т.н., доцент & \makebox[4cm]{\hrulefill} & Ведищев В.В. \\
\end{tabular*}

\vfill
\begin{center}
    Липецк, 2026
\end{center}

\thispagestyle{empty}

\newpage
\tableofcontents
\newpage

\section{Использование метрик качества для текущей оценки достижений}

\subsection{Количественные характеристики качества программного продукта}

Количественные характеристики качества позволяют объективно оценивать состояние программного продукта на различных этапах разработки. В отличие от субъективной экспертной оценки, применение метрик позволяет измерять производительность, стабильность, безопасность и корректность работы системы в числовой форме.

Для системы корпоративного обмена сообщениями количественная оценка качества особенно важна, поскольку продукт должен обеспечивать минимальные задержки передачи данных, устойчивость голосовой связи, корректную доставку сообщений и безопасность обмена информацией внутри корпоративной сети.

Контроль количественных характеристик выполняется на протяжении всего жизненного цикла проекта и используется как для принятия решений о готовности релиза, так и для выявления областей, требующих доработки.

\subsection{Метод «Черного ящика»}

Метод «чёрного ящика» предполагает тестирование системы без анализа её внутренней структуры. Проверке подвергается только внешнее поведение системы: входные данные, реакция системы и соответствие ожидаемому результату.

Для разрабатываемой системы данный метод применяется при тестировании пользовательских сценариев: авторизации, передачи сообщений, голосовой связи и обмена файлами. При этом тестировщик оценивает корректность функционирования системы с позиции конечного пользователя.

Например, выполняется проверка времени доставки текстового сообщения, успешности входа в систему при корректных учётных данных, устойчивости голосового соединения и возможности восстановления связи после сетевого сбоя.

\subsection{Метод «Белого ящика»}

Метод «белого ящика» предполагает тестирование внутренней структуры программной системы, включая исходный код, алгоритмы, ветвления и взаимодействие компонентов.

В рамках проекта метод используется преимущественно бэкенд-разработчиком и специалистом по информационной безопасности для проверки корректности сетевых алгоритмов, обработки ошибок, механизмов аутентификации и устойчивости к отказам.

Тестирование методом белого ящика включает модульное тестирование, анализ покрытия кода, проверку исключительных ситуаций и инспектирование реализации протоколов взаимодействия между компонентами.

\subsection{Примеры метрик качества}

Для оценки качества программного продукта используется набор измеримых метрик, позволяющих определить степень соответствия системы требованиям.

\begin{table}[H]
\centering
\caption{Метрики качества программного продукта}
\begin{tabular}{|p{3.5cm}|p{3.2cm}|p{2.5cm}|p{2.5cm}|p{2.8cm}|}
\hline
\textbf{Проверка} &
\textbf{Метрика} &
\textbf{Хорошо} &
\textbf{Плохо} &
\textbf{Неприемлемо} \\
\hline

Доставка текстового сообщения &
Среднее время доставки &
$<50$ мс &
$50-150$ мс &
$>150$ мс \\
\hline

Голосовая связь &
Средняя задержка аудио &
$<80$ мс &
$80-150$ мс &
$>150$ мс \\
\hline

Передача файлов &
Скорость передачи в LAN &
$>100$ Мбит/с &
$20-100$ Мбит/с &
$<20$ Мбит/с \\
\hline

Авторизация пользователей &
Время входа в систему &
$<1$ с &
$1-3$ с &
$>3$ с \\
\hline

Надёжность доставки сообщений &
Процент успешной доставки &
$>99.9\%$ &
$95-99.9\%$ &
$<95\%$ \\
\hline

Устойчивость системы &
Частота отказов &
$<1$ отказ/месяц &
$1-3$ отказа/месяц &
$>3$ отказов/месяц \\
\hline

Безопасность &
Количество критических уязвимостей &
0 &
1--2 &
$>2$ \\
\hline

Стабильность сборки &
Процент успешных сборок CI &
$>98\%$ &
$90-98\%$ &
$<90\%$ \\
\hline

Покрытие кода тестами &
Процент покрытия &
$>80\%$ &
$50-80\%$ &
$<50\%$ \\
\hline

Использование ресурсов &
Средняя загрузка CPU сервера &
$<50\%$ &
$50-80\%$ &
$>80\%$ \\
\hline

\end{tabular}
\end{table}

Представленные метрики используются для регулярного мониторинга качества системы и позволяют определить степень готовности программного продукта к эксплуатации. Наличие количественных порогов обеспечивает объективность оценки и позволяет классифицировать состояние системы как приемлемое, требующее доработки или критическое.

\end{document}
